\section{Formal Model}\label{sec:s-model}

\subsection{Entries and graphs}

Let $\mathbb E$ be entity identifiers, $\mathbb P$ predicate names, $\mathbb V$ literals (strings,
integers, Booleans, dates, date-times and times in canonical lexical form), $\mathbb W$ message
spans and $\mathbb T$ world times (dates and date-times). An \emph{entry} is one of
\begin{itemize}
\item a claim $c=(s,p,o,\omega,\tau,\vs)$ with $s\in\mathbb E$, $p\in\mathbb P$,
  $o\in\mathbb E\cup\mathbb V$, scope $\omega$ a string or $\bot$ (default scope), valid interval
  $\tau=[\tau^-,\tau^+)$ with $\tau^-\in\mathbb T\cup\{\bot\}$ (unknown) and
  $\tau^+\in\mathbb T\cup\{\bot\}$ (open), both bounds of one datatype;
\item an attribution $a=(u,c,f,\rho,\pi,\gamma,\delta,\vs)$ on claim $c$, with speaker
  $u\in\mathbb E$, function $f$ (request, offer, accept, reject, confirm, inform, commit,
  question), report mode $\rho$ (own, quote, hypothetical), polarity $\pi$ (positive, negative),
  certainty $\gamma$ and derived flag $\delta$;
\item an evidence link $\ell=(c,m,\vs)$ with $m\in\mathbb W$; $\mathrm{id}(\ell)$ is the unified
  identifier of the span (the message, or the message with a character range), so two links can
  share an identifier;
\item a supersession event $x=(c,\kind,\vend,\vs)$ with $\kind\in\{\change,\correction\}$ and
  $\vend\in\mathbb T\cup\{\bot\}$, where $\vend\neq\bot$ only if $\kind=\change$.
\end{itemize}
Every entry $e$ has a version stamp $\vs(e)\ge1$ and an identifier; all fields are fixed when the
entry is written. The stored form also records, for an attribution, the message that expressed
it and, for a supersession event, the replacing claim and the licensing attribution; no query
reads these fields, so they are omitted. A graph $\G$ is a finite set of entries (and of the
entity-equivalence entries of S1.3). For a claim $c$ of $\G$ write
$\att_\G(c)$, $\evs_\G(c)$ and $\sups_\G(c)$ for the attributions, evidence links and
supersession events of $\G$ on $c$, and $N_\G(c)=\{c\}\cup\att_\G(c)\cup\sups_\G(c)$ for its
\emph{neighbourhood}. For a version $r$, $\G{\restriction}r=\{e\in\G:\vs(e)\le r\}$.

\paragraph{Histories.}
A history produces graphs $\G_0=\emptyset\subseteq\G_1\subseteq\dots$, where $\G_v=\G_{v-1}\cup
\Dlt_v$ and every entry of $\Dlt_v$ (the change set of version $v$) has stamp $v$ (assumption
A1). Versions may skip numbers; $\G_v=\G_{v-1}$ then. An entry names only entries of the same or
an earlier change set, so an attribution, link or event on $c$ has a stamp at least $\vs(c)$.
Consequently, for $r_0\le r$,
\begin{equation}\label{eq:truncate}
\G_r{\restriction}r_0=\G_{r_0}.
\end{equation}

\paragraph{Effective valid end.}
The effective valid end $\effend_\G(c)$ of a claim $c$ in $\G$ is $\tau^+(c)$ if
$\tau^+(c)\neq\bot$; otherwise $\vend(x)$ for the first $x\in\sups_\G(c)$ with $\kind(x)=\change$
and $\vend(x)\neq\bot$, in the order of $(\vs(x),\text{identifier})$; and $\bot$ (open, or ended at
an unknown time) if there is no such $x$.
$\effend_\G(c)$ depends only on $N_\G(c)$. Along a history it changes at most once: it is
$\bot$ until the first change event with a stated end and fixed afterwards, because later events
carry larger stamps and entries are never removed (A1 and the order above). Assumption A6 states
that the implementation uses this same order everywhere.

\subsection{Claim patterns and their rows}

A claim pattern is $P=(s_P,p_P,o_P,\omega_P)[u_P,f_P,\rho_P,\pi_P,\delta_P]$ with a label
$\lambda_P$ naming the matched claim. Subject, predicate and object positions are constants or
variables; the scope position is a string, $\mathsf{default}$, the wildcard $\_$ or a variable.
The bracket is optional; if present, at least one of $u_P$ (an entity or a variable) and $f_P$ is
given, and each of $\rho_P,\pi_P,\delta_P$ is either its default ($\own$, positive, false) or
relaxed ($\ast$, any value). A query carries a state $\textsf{st}\in\{\cur,\hist\}$ and an
optional world time $d$, which apply to every pattern.

Let $\equiv$ be an equivalence relation on $\mathbb E$ (S1.3). For a graph $\G$, a \emph{row} of
$P$ is a triple $\mu=(\theta,c,a)$ of a binding $\theta$ of the variables of $P$, a claim
$c\in\G$, and an attribution $a\in\att_\G(c)$ if $P$ has a bracket ($a=\ast$ otherwise), such
that
\begin{description}
\item[(m1) positions.] For each of subject, predicate and object, a constant equals the claim's
  term (entities up to $\equiv$), and a variable $X$ has $\theta(X)$ equal to the claim's term; a
  variable occurring twice therefore forces equality. Scope: a string requires $\omega(c)$ equal
  to it, $\mathsf{default}$ requires $\omega(c)=\bot$, $\_$ accepts any scope, and a variable $X$
  requires $\omega(c)\neq\bot$ and $\theta(X)=\omega(c)$.
\item[(m2) attribution.] If $P$ has a bracket: the speaker of $a$ equals $u_P$ (up to $\equiv$) or
  $\theta(u_P)$ for a variable; $f(a)=f_P$ if $f_P$ is given; and each of $\rho_P,\pi_P,\delta_P$
  that is not relaxed equals the corresponding field of $a$.
\item[(m3) state.] If $\textsf{st}=\cur$, $\sups_\G(c)=\emptyset$; if $\textsf{st}=\hist$, no
  $x\in\sups_\G(c)$ has $\kind(x)=\correction$.
\item[(m4) world time.] If $d$ is given, $c$ is not known to lie outside its interval at $d$:
  not $\tau^-(c)>d$, not $\tau^+(c)\le d$, and if $\tau^+(c)=\bot$, not $\effend_\G(c)\le d$.
  A bound that is $\bot$, or a comparison between a date and a date-time, excludes nothing.
\end{description}
$\rows{P}{\G}{}$ is the set of rows. We say $c$ \emph{matches} $P$ in $\G$, written
$c\matches_\G P$, if some $\theta$ and $a$ satisfy (m1) and (m2); (m3), (m4) and filters are
ignored. Note that (m1)--(m4) refer only to $N_\G(c)$ and to $\equiv$ (Lemma~\ref{lem:locality}).

\subsection{Entity equivalence and vocabulary}

Entity merges are recorded as equivalence entries $(e_1,e_2,\vs)$; $\equiv_\G$ is the equivalence
closure of those in $\G$. Constants match every member of their class, and an entity answer is
reported as the least member of its class. The vocabulary (predicate names, aliases,
exclusivity rules) is a version label of the history. Equivalence entries and vocabulary changes
are not among the four kinds of entries: the proofs assume that neither occurs between the two
versions compared (assumption A5), and the system invalidates every view when one occurs.

\subsection{Queries}

A query of the fragment has a body
\[
B=\sigma_\Phi(P_1\bowtie\dots\bowtie P_k)\;\rhd_\Psi\;N,\qquad 1\le k\le 3,
\]
where $P_1,\dots,P_k$ are the positive patterns, $\Phi$ a conjunction of filters, and the optional
anti-join $\rhd_\Psi N$ is the single \textsc{not exists} over pattern $N$ with its own filters
$\Psi$; $N$ shares at least one variable with the positive part. A \emph{filter} compares a
variable with a constant ($=,\neq$ for entities, up to $\equiv$; $=,\neq,<,\le,>,\ge$ for literals
of one datatype), or a valid-time bound $\mathrm{start}(\lambda)$ or $\mathrm{end}(\lambda)$ of a
labelled claim with a date or date-time constant. $\mathrm{start}(\lambda)\mathrel{op}d$ holds iff
$\tau^-(c)\neq\bot$ and $\tau^-(c)\mathrel{op}d$. $\mathrm{end}(\lambda)\mathrel{op}d$ uses
$\tau^+(c)$ if set, else $\effend_\G(c)$ if set, else the open end, which is later than every
time ($>,\ge,\neq$ hold, $<,\le,=$ do not). A comparison that is undefined (mixed datatypes) is
false. No filter refers to the time of reading.

The rows of the body are tuples $\vec\mu=(\mu_1,\dots,\mu_k)$ with $\mu_i\in\rows{P_i}{\G}{}$
whose bindings agree on shared variables, that satisfy $\Phi$, and for which no row
$\nu\in\rows{N}{\G}{}$ agrees with $\vec\mu$ on shared variables and satisfies $\Psi$ together
with $\vec\mu$; variables only in $N$ are existential. Write $\rows{B}{\G}{}$ for this set,
$\theta(\vec\mu)$ for the joint binding, and
\[
W(\vec\mu)=\{c_1,\dots,c_k\}\cup\{a_i:a_i\neq\ast\}
\]
for its witness. The root of the query is one of
\begin{itemize}
\item[] Before the root operator, every entity value bound to an answer variable or to the counted
  variable is replaced by the representative of its class (the least member), so grouping,
  minimisation and distinct counting see representatives; $\theta(\vec\mu)$ below denotes the
  binding after this replacement. Join variables still compare entities as written.
\item a projection on variables $X$: the answers are $\{\theta(\vec\mu)|_X\}$, and the witnesses
  of answer $t$ are the $\subseteq$-minimal sets among $\{W(\vec\mu):\theta(\vec\mu)|_X=t\}$;
\item a distinct count of variable $y$ grouped by $X$: for each group $g$, with members
  $M_g=\{\theta(\vec\mu)(y):\theta(\vec\mu)|_X=g\}$, the answer is $(g,|M_g|)$ and its single
  witness is the union over $m\in M_g$ of the minimal witnesses of the rows with
  $\theta|_X=g$ and $\theta(y)=m$; without grouping and without rows the answer is $(0)$ with no
  witness;
\item a latest-state selection on the claim labelled $\lambda_j$, followed by a projection. The
  key of a row is $\max(\vs(c_j),\vs(a_j))$ for selection by version (the version at which the
  pattern became matched, so a later acceptance of an older offer counts at the acceptance), or
  $\tau^-(c_j)$ for selection by valid time, where rows with $\tau^-(c_j)=\bot$ are dropped. The
  key of a claim is the maximum key of its rows; the rows kept are all rows of the claims with
  maximum key (ties keep every tied claim), and the projection then applies to them.
\end{itemize}
The answer set is
$\ans(\Q,\G)$ and $\wit(\Q,\G)$ is the set of pairs (answer, its witnesses). The evidence is recorded per answer: for an answer $t$,
\[
\ev_t(\Q,\G)=\{\mathrm{id}(\ell):\ell\in\evs_\G(c),\ c\text{ a claim of a witness of }t\},
\]
and $\ev(\Q,\G)$ is the set of pairs $(t,\ev_t(\Q,\G))$. Recording evidence per answer, not as
one set over all answers, makes a new citation of a span visible even when the span is already
cited for another answer, as the sentence-level citations of a view require.
Attributions in witnesses contribute no evidence of their own. A query written
\texttt{WITHOUT PROVENANCE} has $\ev_t=\emptyset$ for every answer by definition; its witnesses are still computed.

\paragraph{Cut-off.}
$\cut{\Q}{r}$ evaluated on $\G$ is, by definition, $\Q$ evaluated on $\G{\restriction}r$ (with
$\equiv_{\G{\restriction}r}$). The implementation realises this by comparing the stamp of every
entry the evaluation touches, including supersession events, evidence links and equivalence
entries, with the cut-off placeholder (assumption A2).

\paragraph{Local filters.}
A filter is \emph{local} to pattern $P$ if its only variable or label belongs to $P$. Since
$\sigma_\varphi(R\bowtie S)=\sigma_\varphi(R)\bowtie S$ when $\varphi$ mentions only variables of
$R$, local filters may be moved into their pattern; write $P^\Phi$ for $P$ with its local filters,
and $\rows{P^\Phi}{\G}{}$ for the rows of $P$ that satisfy them.
